% Generated draft: requires mathematical and source review.
\documentclass[11pt]{article}
\usepackage[margin=25mm]{geometry}
\usepackage{fontspec}
\setmainfont{Latin Modern Roman}
\usepackage{amsmath,amssymb,mathtools,unicode-math}
\setmathfont{Latin Modern Math}
\usepackage{enumitem,needspace,xurl,hyperref}
\setlength{\parindent}{0pt}
\setlength{\parskip}{5pt}
\setlength{\emergencystretch}{3em}
\newcommand{\N}{\mathbb{N}}
\newcommand{\Z}{\mathbb{Z}}
\newcommand{\Q}{\mathbb{Q}}
\newcommand{\R}{\mathbb{R}}
\newcommand{\C}{\mathbb{C}}
\newcommand{\F}{\mathbb{F}}
\newcommand{\A}{\mathbb{A}}
\newcommand{\PP}{\mathbb{P}}
\newcommand{\E}{\mathbb{E}}
\newcommand{\eps}{\varepsilon}
\newcommand{\dd}{\,\mathrm{d}}
\begin{document}
\section*{UnsolvedMath ID: 20000451}
\textbf{Source problem code:} catalogue entry
\subsection{Definitions and mathematical statement}
For a separable polynomial \(g(t)\in\Q[t]\) of degree at least three, let \(C\) be the smooth projective hyperelliptic curve associated to \(s^2=g(t)\), and let \(\operatorname{Jac}(C)\) be its Jacobian. An isogeny over \(\Q\) is a surjective homomorphism of abelian varieties over \(\Q\) with finite kernel. Silverberg's question asks for such a \(g\), an elliptic curve \(E/\Q\), an abelian variety \(B/\Q\), and an isogeny
\[
 \operatorname{Jac}(C)\sim_{\Q}E^r\times B,
 \qquad r\ge4\quad\text{or}\quad(r=3\text{ and }\dim B\le1).
\]
The trivial abelian variety is allowed when \(\dim B=0\). There is no dimension bound on \(B\) in the \(r\ge4\) alternative.

Ellenberg's related question asks for a map \(\PP^1\to E^4/\{\pm1\}\), with simultaneous inversion on all four factors, whose image is contained in no \(H/\{\pm1\}\) for a proper subgroup \(H\subsetneq E^4\). For precision a geometric version should specify a nonconstant morphism, algebraic subgroups, and the ground field/field of definition of \(H\). Those conventions are not explicit in the source; the surrounding rational construction suggests \(\Q\), while a separate remark explicitly asks about \(\C\). The later report claims a Jacobian--Kummer equivalence between a surjection \(\operatorname{Jac}(C)\to E^r\) and a nondegenerate rational curve on \(E^r/\{\pm1\}\), under nonconstancy and field conventions not supplied here. No unconditional equivalence with all descent and basepoint issues suppressed is asserted in this draft.

The historical remarks report that \((r,\dim B)=(2,0)\) and \((3,2)\) are possible; that such constructions would improve lower bounds for the density of quadratic twists of \(E\) of rank at least \(r\); and that over \(\C\), \((3,0)\) is possible, while large \(r\) might be impossible. They suggest composite-exponent curves \(s^2=t^k+1\), including the map toward \(s^2=t^3+1\) when \(3\mid k\), and insist that hyperellipticity is needed for the applications. A random-matrix heuristic is said to suggest a positive-power density bound, without a specified exponent. Another proposed route is to start with pairwise \(\Q\)-isogenous elliptic curves \(E_1,E_2,E_3\), seek a nonsplit principal polarization on \(A=E_1\times E_2\times E_3\), and then seek a genus-three hyperelliptic Jacobian realization. This is a proposed route, not a verified descent or hyperellipticity theorem.

The supplied partial-result report analyzes
\[
 C_0:y^2=x^8+14x^4+1,\quad
 E:Y^2=X^3-5616X+120960,\quad
 E^{(-1)}:Y^2=X^3-5616X-120960.
\]
It claims the exact arithmetic decomposition
\[
 \operatorname{Jac}(C_0)\sim_{\Q}E^2\times E^{(-1)},
\]
from three \(\Q\)-defined involution quotients. At the common good prime eleven it gives Frobenius traces \(a_{11}(E)=4\), \(a_{11}(E^{(-1)})=-4\), where \(a_p(E)=p+1-\#E(\F_p)\), certifying that these factors are not \(\Q\)-isogenous. Thus this geometric elliptic-cube example is not a \(\Q\)-elliptic cube. A separate prime-three trace obstruction is reported for a standard genus-five geometric fifth-power example and a displayed \(E_5^5\), but neither that curve nor \(E_5\) appears in this record. These descent obstructions do not supply the requested \(\Q\)-example or rule out all others.
\subsection{Short English statement}
Find a hyperelliptic Jacobian over \(\Q\) with a large repeated elliptic factor, or a related nondegenerate rational curve on a Kummer quotient. The supplied report shows that a familiar geometric elliptic cube decomposes into nonisogenous quadratic-twist factors over \(\Q\), illustrating a descent obstruction rather than solving the construction problem.
\subsection{Sources}
\begin{itemize}
\item Local UnsolvedMath registry, permanent ID \(20000451\), source code \texttt{catalogue entry}; AIM algebraic number theory corpus, \texttt{aim-algebraic-number-theory-notes.json}, notes[102]. Attribution: Ulam AI / UnsolvedMath Contributors, CC BY 4.0.
\item AIM \emph{Rational and integral points on higher dimensional varieties}, Question 18(a) attributed to Silverberg and (b) to Ellenberg: \url{https://aimath.org/WWN/qptsurface2/qptsurface2.pdf}. Remarks name Silverberg, Voloch, Poonen, Ellenberg and Rodriguez-Villegas. The verification reports recovery of superscripts and strict containment from HTML/PDF; no new inspection is claimed. The remarks contain page artifacts and end with ``Problem/''.
\item Original literature field is empty. The report supplies no primary references, quotient formulas, genus-five equation or \(E_5\) definition for its additional claims.
\item Registry dated August 14, 2026: statement exact, status \texttt{partially\_solved}, classification \texttt{PARTIAL-PROGRESS}; attempt 1 records problem status open, partial result, rigor proved, novelty confidence low. Its decomposition, trace certificates and Jacobian--Kummer criterion are machine-generated research aids, not peer reviewed, requiring independent verification before citation. The report's statement that its obstruction was not located in the literature checked is a supplied novelty assessment only.
\end{itemize}
\end{document}
